\section{数学规划与混合整数规划基础}

数学规划是运筹学的核心工具，用于在满足一组约束条件的前提下，寻找使目标函数取得最优值的决策方案。一个数学规划模型由三个基本要素构成：决策变量，描述决策者可以控制的量；目标函数，衡量决策方案优劣的准则；约束条件，刻画决策必须满足的限制。

根据决策变量的类型和约束的结构，数学规划可分为若干类别。线性规划（LP）要求目标函数和约束均为决策变量的线性函数，且变量取值连续\citep{bertsimas1997introduction}。整数规划（IP）在线性规划的基础上要求部分或全部变量取整数值，当所有变量均为整数时称为纯整数规划。混合整数规划（MIP）则同时包含连续变量和整数变量，是实际应用中最常见的形式\citep{wolsey1998integer, nemhauser1988integer}。一般的混合整数规划可以写成如下形式：
\begin{equation}
\begin{aligned}
\min \quad & c^\top x + d^\top y \\
\text{s.t.}\quad & Ax + By \ge b,\\
& x \ge 0,\\
& y \in \{0,1\}^n.
\end{aligned}
\label{eq:general_mip}
\end{equation}
其中$x$表示连续决策变量，$y$表示二元或整数决策变量，目标函数表示成本、收益或两者的组合，约束表示资源、容量、逻辑关系等限制。

为帮助理解 MIP 的结构，考虑一个固定费用网络设计问题的简化形式：
\begin{equation}
\begin{aligned}
\min \quad & \sum_{(i,j)\in A} F_{ij} y_{ij} + \sum_{(i,j)\in A} c_{ij} f_{ij} \\
\text{s.t.}\quad & \sum_{j:(i,j)\in A} f_{ij} - \sum_{j:(j,i)\in A} f_{ji} = b_i, \quad \forall i\in N,\\
& f_{ij} \le U_{ij}\, y_{ij}, \quad \forall (i,j)\in A,\\
& f_{ij}\ge 0,\quad y_{ij}\in\{0,1\}.
\end{aligned}
\label{eq:fcnd}
\end{equation}
其中$y_{ij}=1$表示弧$(i,j)$被开通，$f_{ij}$表示弧上的流量。约束$f_{ij}\le U_{ij}y_{ij}$是典型的连续流量变量与二元开通变量的耦合：只有弧被开通时流量才能为正。本文中的服务弧激活变量$x$与顾客选择份额$w$的耦合，本质上也属于这类 MIP 结构，只是更为复杂——因为需求不是外生给定的，而是由顾客选择内生产生的。例如，后文模型中的容量约束将顾客选择份额所诱导的客流量限制在已激活服务弧的座位容量之内，其作用与固定费用网络设计中的$f_{ij}\le U_{ij}y_{ij}$类似，都是通过整数变量控制连续流量变量的可行范围。

\section{网络流与服务网络设计}

网络流模型是描述资源在网络中流动的基本框架\citep{ahuja1993network}。一个网络由节点集合$N$和弧集合$A$构成，流在弧上传输并在节点处满足守恒条件。最小费用网络流问题的标准形式为：
\begin{equation}
\begin{aligned}
\min \quad & \sum_{(i,j)\in A} c_{ij} f_{ij}\\
\text{s.t.}\quad & \sum_{j:(i,j)\in A} f_{ij} - \sum_{j:(j,i)\in A} f_{ji} = b_i, \quad \forall i\in N,\\
& 0\le f_{ij}\le u_{ij}, \quad \forall (i,j)\in A.
\end{aligned}
\label{eq:mcf_basic}
\end{equation}
其中$b_i$为节点$i$的净供给量（正为供给，负为需求），$u_{ij}$为弧容量。当网络中存在多种商品（即多个起讫对的流量需求）时，每个商品$k$各自满足流守恒但共享弧容量，问题扩展为多商品流（MCF）问题。

服务网络设计（SND）在多商品流的基础上引入了服务选择决策\citep{magnanti1984network, crainic2000service}。与经典网络流问题中所有弧均可用不同，SND中的服务弧需要通过分配资源来激活——只有被激活的服务弧才能承载流量。因此，SND不仅要确定流量的路径分配，还要决定哪些服务弧被激活、以何种频率运营、分配何种类型的资源。这使得SND成为一个同时包含网络设计决策和流量分配决策的优化问题。

在本文的问题中，网络节点对应不同时间点的服务站点，服务弧对应计划中的服务班次，资源对应可分配的服务设备。运营商需要在资源总量、流守恒和运营约束下，确定最优的服务激活和资源分配方案。

从问题结构的递进关系来看，多商品流问题假设网络中所有弧均可用，只需确定流量分配；服务网络设计在此基础上引入服务激活决策，需要同时确定哪些弧可用以及流量如何分配；而考虑顾客选择的服务网络设计则进一步将需求从外生输入变为内生变量，使服务激活、资源分配和需求分配三者相互耦合。因此，从网络流到服务网络设计，再到考虑顾客选择的服务网络设计，问题复杂度的增加主要来自两个方面：一是服务开通决策引入整数变量；二是顾客选择行为使需求分配由外生参数变为内生变量。

\section{时空网络建模}

交通服务具有天然的时间维度：同一条线路在不同时间出发构成不同的服务选项。为在优化模型中精确表示这种时间结构，时空网络建模方法将规划时域离散化为若干时间点，并将空间网络在时间维度上展开。

具体而言，对于空间网络中的每个节点$i$和离散化后的每个时间点$t$，构造一个时空节点$(i,t)$。网络中的弧分为两类：服务弧连接不同空间节点在不同时间点的时空节点，表示从$(i,t)$出发、在$(j,t')$到达的一次服务运营，其中$t'-t$等于该服务的旅行时间；持有弧连接同一空间节点在相邻时间点的时空节点，表示资源在该节点的等待。如图~\ref{fig:time_space_network_basic}所示，空间节点$i$和$j$在不同时间点被复制为多个时空节点；从$(i,t_1)$到$(j,t_3)$的斜向弧表示一次服务运营，而同一空间节点相邻时间点之间的水平弧表示资源等待。

\begin{figure}[H]
\centering
\begin{tikzpicture}[scale=0.95,
    timepoint/.style={circle, fill=black, inner sep=1.6pt},
    service/.style={->, thick, blue!70!black},
    hold/.style={->, thick, gray!70},
    labelnode/.style={font=\small}]
% node labels
\node[labelnode] at (-0.6,1.6) {节点 $j$};
\node[labelnode] at (-0.6,0) {节点 $i$};
% time labels
\foreach \x/\t in {0/$t_1$,1.6/$t_2$,3.2/$t_3$,4.8/$t_4$,6.4/$t_5$} {
    \node[labelnode] at (\x,-0.45) {\t};
}
% time-space nodes
\foreach \x in {0,1.6,3.2,4.8,6.4} {
    \node[timepoint] at (\x,0) {};
    \node[timepoint] at (\x,1.6) {};
}
% holding arcs
\foreach \x/\xn in {0/1.6,1.6/3.2,3.2/4.8,4.8/6.4} {
    \draw[hold] (\x,0) -- (\xn,0);
    \draw[hold] (\x,1.6) -- (\xn,1.6);
}
% service arcs
\draw[service] (0,0) -- (3.2,1.6);
\draw[service] (1.6,0) -- (4.8,1.6);
\draw[service] (3.2,0) -- (6.4,1.6);
% annotations
\node[labelnode, blue!70!black] at (2.25,1.15) {服务弧};
\node[labelnode, gray!70!black] at (5.5,0.25) {持有弧};
\draw[->, blue!70!black, thin] (2.25,1.0) -- (1.5,0.75);
\draw[->, gray!70!black, thin] (5.5,0.15) -- (5.6,0.0);
\node[labelnode] at (3.2,-0.9) {时间};
\end{tikzpicture}
\caption{时空网络示意图}
\label{fig:time_space_network_basic}
\end{figure}

时空网络的优势在于将时间约束（如最小间隔时间、中转衔接时间）转化为网络结构约束，使得复杂的调度问题可以在统一的网络流框架下建模\citep{desrosiers1995time, zhu2014scheduled}。然而，时空网络的规模随时间分辨率的提高而快速增长：若空间网络有$n$个节点、离散化产生$T$个时间点，则时空网络包含$nT$个节点，弧的数量也相应增加。这种规模膨胀是本文需要解决的关键计算挑战之一\citep{boland2019price}。

\section{离散选择模型基础}

在经典的服务网络设计模型中，需求被视为外生给定的——每个起讫对的客流量固定，不受服务计划的影响。然而在客运场景中，需求是内生的：顾客会根据可用行程的服务属性进行选择。离散选择模型为描述这种选择行为提供了理论基础。

离散选择模型假设每位顾客面对一个有限的选择集合，集合中的每个方案具有一个效用值，由系统效用（可观测属性的函数）和随机误差项组成\citep{mcfadden1974conditional, benakiva1985discrete, train2009discrete}。顾客选择效用最高的方案。当随机误差项服从独立同分布的极值分布时，得到多项Logit（MNL）模型，其选择概率具有封闭的解析表达式：方案$i$被选择的概率正比于其吸引力$\alpha_i = \exp(V_i)$，其中$V_i$为系统效用。

在本文的问题中，每个市场中的顾客面对若干可行行程和一个外部选项（不出行）。外部选项的吸引力反映了顾客放弃出行的倾向。为了将选择行为嵌入线性优化模型，网络收入管理领域发展了基于选择的线性规划\citep{liu2008choice}和基于销售的线性规划（SBLP）\citep{gallego2015general}等框架，通过吸引力参数将非线性选择概率转化为线性约束。后续模型采用基于吸引力的线性表示，通过吸引力参数和选择一致性约束刻画顾客的需求分配行为。

\section{松弛、上下界与最优性间隙}

在混合整数规划求解中，直接获得全局最优解往往十分困难。为衡量一个可行解的质量，通常使用上界、下界和最优性间隙等概念。

对于最小化问题而言，任意可行解的目标函数值均可作为原问题最优值的上界；而若能够证明原问题的最优值不可能低于某一数值，则该数值称为下界。当上界与下界足够接近时，即可说明当前可行解已经接近最优。

松弛模型是构造下界的常用方法。其基本思想是放宽原问题中的部分约束，得到一个更容易求解的问题。由于松弛模型的可行域包含原问题的可行域，在最小化问题中，松弛模型的最优值不高于原问题最优值，因此可作为有效下界。另一方面，上界必须由原问题中的实际可行解给出。

本文后续提出的求解框架正是基于这一思想：首先在聚合网络和超级行程空间中求解松弛模型以获得下界；然后在完整网络上修复得到可行解以获得上界；最后通过上下界之间的差距判断当前解是否达到预设最优性容差。

在实际求解混合整数规划时，分支定界（Branch-and-Bound）是最常用的精确方法\citep{land1960automatic}。该方法通过系统地将问题分解为子问题（分支），并利用松弛模型的最优值对子问题进行剪枝（定界），逐步缩小搜索空间。在求解过程中，算法维护两个关键量：当前最好可行解的目标值（incumbent，即上界）和所有未探索子问题松弛值中的最小值（即下界）。二者之差的相对值称为最优性间隙（optimality gap），用于衡量当前解距离全局最优的程度。当间隙降至预设容差以下时，即可认为当前解已达到实际可接受的最优性水平。

当问题具有特殊的分块结构时，列生成（Column Generation）是另一种重要的精确求解技术\citep{desaulniers2005column, barnhart1998branch}。列生成通过将大规模线性规划分解为限制主问题和定价子问题，避免显式枚举所有变量，仅在需要时动态生成有改进潜力的列。当列生成与分支定界结合时，称为分支定价（Branch-and-Price），适用于变量数量极大的整数规划问题。

对于本文研究的大规模混合整数规划，由于行程变量数量庞大且线性松弛较弱，商业求解器的分支定界过程往往进展缓慢：每个节点的求解代价高，且松弛界改善有限。因此，本文不依赖通用分支定界直接求解完整模型，而是通过超级行程聚合构造更紧的松弛下界，并通过频率固定和收入约束在完整网络上高效生成可行上界，从而在框架层面实现间隙的快速收敛。
